% Group 4 section for the combined CNS phishing-detection paper.
% Self-contained IEEE conference document. Replace the figure paths with the
% generated figure files when they are available in the Overleaf project.
% All reported experimental values below are transcribed from the supplied
% Group 4 experiment specification; they are not independently recomputed here.

\documentclass[conference]{IEEEtran}
\usepackage[T1]{fontenc}
\usepackage{amsmath,amssymb}
\usepackage{array,booktabs}
\usepackage{graphicx}
\usepackage{url}
\usepackage{cite}

\title{Group 4: Linear Discriminant Analysis for Low-Overhead Real-Time Web Filter Applications}
\author{\IEEEauthorblockN{Group 4}\IEEEauthorblockA{Bharatiya Vidya Bhavan's Sardar Patel Institute of Technology, Mumbai}}

\begin{document}
\maketitle

\subsection{Group 4: Linear Discriminant Analysis for Low-Overhead Real-Time Web Filter Applications}
\label{subsec:group4-lda}

\subsubsection{Motivation and Problem Formulation}
Phishing website detection is formulated as binary classification of a website record into a legitimate class or a phishing class. Linear Discriminant Analysis (LDA) is considered as a compact statistical classifier for settings in which an inexpensive decision stage is desirable, including browser-side filters, DNS security gateways, and edge systems. Unlike an ensemble of decision trees, a fitted LDA prediction can be expressed as a linear score over the supplied feature vector. This simple scoring stage motivates evaluation as a possible first-line filter; it does not, by itself, establish end-to-end wire-speed performance because feature extraction and system-level processing also consume resources.

The assigned method is classification LDA. It should be distinguished from using LDA solely as a dimensionality-reduction transform before a separate classifier. Here, LDA estimates class-conditional distributions and directly assigns a class (and, when requested, posterior estimates). The model assumes class-conditional multivariate Gaussian distributions with a common covariance matrix, which produces a linear decision boundary.

\subsubsection{Dataset and Leakage-Safe Preprocessing}
Experiments use the Kaggle \emph{Phishing Website Detector} dataset in \texttt{phishing.csv}. The data contain 11,054 samples and 32 original columns: 30 extracted tabular website indicators, a sequential \texttt{Index} identifier, and the target. The indicators span URL structure, domain and website properties, security-related signals, and external reputation or traffic information. The \texttt{Index} field has 11,054 unique sequential values and is excluded from the predictors as a non-predictive row identifier.

The ground-truth labels are mapped as follows: original label $-1$ (phishing) becomes positive class $y=1$, and original label $+1$ (legitimate) becomes class $y=0$. There are 4,897 phishing samples (44.3\%) and 6,157 legitimate samples (55.7\%). A stratified 80/20 train--test split with \texttt{random\_state=42} gives 8,843 training samples and 2,211 held-out test samples. Standardization is encapsulated in a scikit-learn \texttt{Pipeline} and fitted only on training data, including within each cross-validation fold. Thus, the held-out test partition is not used to estimate preprocessing statistics or select the solver.

\subsubsection{Mathematical Formulation}
Let $\mathbf{x}_i\in\mathbb{R}^{p}$ denote the feature vector for observation $i$, where $p=30$, and let $y_i\in\{0,1\}$ denote its class, with $1$ representing phishing. Let $n_k$ and $\boldsymbol{\mu}_k$ be the number and mean vector of samples in class $k$, respectively, and let $\boldsymbol{\mu}$ denote the overall mean. The within-class scatter matrix is
\begin{equation}
 S_W = \sum_{k\in\{0,1\}}\sum_{i:y_i=k}
 (\mathbf{x}_i-\boldsymbol{\mu}_k)(\mathbf{x}_i-\boldsymbol{\mu}_k)^T,
 \label{eq:within-scatter}
\end{equation}
which measures the aggregate dispersion of samples around their own class means. The between-class scatter matrix is
\begin{equation}
 S_B = \sum_{k\in\{0,1\}} n_k
 (\boldsymbol{\mu}_k-\boldsymbol{\mu})
 (\boldsymbol{\mu}_k-\boldsymbol{\mu})^T,
 \label{eq:between-scatter}
\end{equation}
which measures class-centroid separation relative to the overall mean.

Fisher's discriminant criterion for a projection vector $\mathbf{w}$ is
\begin{equation}
 J(\mathbf{w}) =
 \frac{\mathbf{w}^T S_B\mathbf{w}}{\mathbf{w}^T S_W\mathbf{w}}.
 \label{eq:fisher-criterion}
\end{equation}
The stationary directions satisfy the generalized eigenvalue equation
\begin{equation}
 S_B\mathbf{w}=\lambda S_W\mathbf{w},
 \label{eq:generalized-eigenproblem}
\end{equation}
where $\lambda$ is an eigenvalue. For two classes and nonsingular $S_W$, the discriminant direction is proportional to
\begin{equation}
 \mathbf{w}\propto S_W^{-1}
 (\boldsymbol{\mu}_1-\boldsymbol{\mu}_0).
 \label{eq:two-class-direction}
\end{equation}
In binary classification, $\operatorname{rank}(S_B)\leq 1$, so there is at most one nonzero discriminant direction (exactly one only when the class means differ in a direction retained by the data). The SVD solver used here is rank-aware and avoids explicitly forming the matrix inverse in (\ref{eq:two-class-direction}).

Under class-conditional Gaussian distributions with shared covariance $\boldsymbol{\Sigma}$, the class-$k$ discriminant score is
\begin{equation}
 \delta_k(\mathbf{x}) = \mathbf{x}^T\boldsymbol{\Sigma}^{-1}\boldsymbol{\mu}_k
 -\frac{1}{2}\boldsymbol{\mu}_k^T\boldsymbol{\Sigma}^{-1}\boldsymbol{\mu}_k
 +\ln(\pi_k),
 \label{eq:lda-score}
\end{equation}
where $\pi_k$ is the prior probability of class $k$. The predicted class is the one with the larger score. In the log-posterior comparison, the shared quadratic term $-\tfrac{1}{2}\mathbf{x}^T\boldsymbol{\Sigma}^{-1}\mathbf{x}$ appears in both classes and cancels. The remaining score difference is affine in $\mathbf{x}$, hence the decision boundary is linear.

For the phishing-positive convention, define $\mathbf{w}=\boldsymbol{\Sigma}^{-1}(\boldsymbol{\mu}_1-\boldsymbol{\mu}_0)$ and write the binary score as $g(\mathbf{x})=\mathbf{w}^T\mathbf{x}+b$. The reported intercept is
\begin{equation}
 b=-\mathbf{w}^T\left(\frac{\boldsymbol{\mu}_1+\boldsymbol{\mu}_0}{2}\right)
 +\ln\left(\frac{\pi_1}{\pi_0}\right)=-0.7779.
 \label{eq:lda-intercept}
\end{equation}
The sign convention is such that a positive score favors the phishing class. The numeric intercept is reported as supplied for this experiment.

\subsubsection{Solver Selection}
Two solver configurations were compared using five-fold stratified cross-validation on the training partition only. The SVD solver is rank-aware and does not require an explicit covariance inverse; the LSQR solver was evaluated with Ledoit--Wolf automatic shrinkage (\texttt{shrinkage=auto}) to regularize covariance estimation. Table~\ref{tab:solver-cv} reports the supplied mean F1 and fold standard deviation. The SVD configuration was selected because it had the slightly higher mean cross-validation F1.

\begin{table}[!t]
\caption{Training-Only Stratified Cross-Validation Solver Comparison}
\label{tab:solver-cv}
\centering
\footnotesize
\begin{tabular}{lcc}
\toprule
\textbf{LDA configuration} & \textbf{Mean CV F1} & \textbf{Fold standard deviation}\\
\midrule
\texttt{solver=svd} & 0.9099 & 0.0051\\
\texttt{solver=lsqr}, \texttt{shrinkage=auto} & 0.9098 & 0.0052\\
\bottomrule
\end{tabular}
\end{table}

\subsubsection{Held-Out Test Results}
After solver selection, the chosen pipeline was evaluated on the held-out test set of $N=2{,}211$ samples. Table~\ref{tab:lda-performance} gives the supplied metrics. Phishing is the positive class throughout.

\begin{table}[!t]
\caption{Performance on the Held-Out Test Set ($N=2{,}211$)}
\label{tab:lda-performance}
\centering
\footnotesize
\begin{tabular}{lc}
\toprule
\textbf{Metric} & \textbf{Value}\\
\midrule
Accuracy & 91.68\% (0.9168)\\
Precision & 92.42\% (0.9242)\\
Recall (sensitivity) & 88.46\% (0.8846)\\
F1-score & 90.40\% (0.9040)\\
Specificity & 94.24\% (0.9424)\\
ROC-AUC & 0.9738\\
PR-AUC & 0.9694\\
\bottomrule
\end{tabular}
\end{table}

The confusion matrix in Table~\ref{tab:lda-confusion} shows 113 false negatives (phishing websites classified as legitimate) and 71 false positives (legitimate websites classified as phishing). False negatives can expose users to credential theft or other malicious activity; false positives can block legitimate services and reduce user trust. Both error types matter, and the acceptable operating point depends on the deployment's risk tolerance.

\begin{table}[!t]
\caption{Held-Out Test Confusion Matrix}
\label{tab:lda-confusion}
\centering
\footnotesize
\begin{tabular}{lcc}
\toprule
\textbf{Actual class} & \textbf{Predicted legitimate (0)} & \textbf{Predicted phishing (1)}\\
\midrule
Legitimate (0) & 1,161 (TN) & 71 (FP)\\
Phishing (1) & 113 (FN) & 866 (TP)\\
\bottomrule
\end{tabular}
\end{table}

\subsubsection{Computational Efficiency}
The supplied local benchmark values are summarized in Table~\ref{tab:lda-efficiency}. The per-sample mean is reported as 0.0064~ms (6.37~$\mu$s), while batch inference over 2,211 samples took 0.0141~s. These values describe the reported model inference measurement; hardware details and upstream feature-extraction cost were not included in the supplied benchmark specification and therefore are not inferred here.

\begin{table}[!t]
\caption{Computational Efficiency Measurements}
\label{tab:lda-efficiency}
\centering
\footnotesize
\begin{tabular}{lc}
\toprule
\textbf{Parameter} & \textbf{Reported value}\\
\midrule
Total training time & 0.0635~s ($\sim$64~ms)\\
Test inference time ($2{,}211$-sample batch) & 0.0141~s\\
Mean latency per sample & 0.0064~ms (6.37~$\mu$s)\\
Throughput & 156,898 samples/s\\
\bottomrule
\end{tabular}
\end{table}

The single linear score requires a weighted sum over the input features after those features have been obtained. By contrast, tree ensembles evaluate multiple branch paths across many component trees, which introduces a different inference structure and typically requires storing many tree nodes. This structural difference makes LDA a plausible lightweight first-stage filter in a cascade with a more complex classifier for uncertain cases. A direct runtime or memory comparison against Gradient Boosting and Extra Trees was not supplied, so no measured superiority over those cohort models is claimed. Likewise, the reported microsecond-scale model latency should not be interpreted as a demonstrated packet-inspection rate without deployment hardware, parser cost, request handling, and concurrency measurements.

\subsubsection{Discriminant Coefficients}
Table~\ref{tab:lda-coefficients} lists selected supplied coefficients. Their signs and magnitudes are conditional on the dataset's feature encoding, preprocessing, class order, and fitted model. They describe the model's linear score and should not be interpreted as causal effects.

\begin{table*}[!t]
\caption{Selected Linear Discriminant Coefficients and Security Interpretation}
\label{tab:lda-coefficients}
\centering
\footnotesize
\begin{tabular}{lcl}
\toprule
\textbf{Feature} & \textbf{Coefficient $w_j$} & \textbf{Security interpretation (as supplied)}\\
\midrule
\texttt{HTTPS} & $-2.6774$ & Strongest indicator of legitimate SSL trust\\
\texttt{AnchorURL} & $-2.3623$ & Internal domain anchoring indicates legitimacy\\
\texttt{ServerFormHandler} & $-0.6955$ & External or blank action URLs indicate data exfiltration\\
\texttt{ShortURL} & $+0.6597$ & Obfuscated shortening links strongly indicate phishing\\
\texttt{WebsiteTraffic} & $-0.6143$ & High Alexa/traffic rank correlates with legitimate sites\\
\bottomrule
\end{tabular}
\end{table*}

\subsubsection{Operational Considerations and Limitations}
The default decision boundary corresponds to a score or posterior threshold determined by the fitted classifier. An operator could consider lowering a phishing-probability threshold from 0.5 to 0.3 when prioritizing recall. In general, lowering the threshold tends to catch more phishing samples while also risking more false positives; the resulting recall, precision, and specificity must be measured on validation data at the chosen threshold before operational use. The supplied test metrics correspond to the reported experiment and do not establish the performance at a 0.3 threshold.

LDA's linear scoring stage is computationally compact relative to multi-tree ensembles and can be considered as an early filter for resource-constrained settings. However, this study's measurements do not include a controlled head-to-head latency or memory benchmark of the cohort algorithms, nor do they establish wire-speed processing on network hardware. The low-latency claim should therefore be restricted to the supplied model-inference measurement until end-to-end tests are conducted on the target browser, DNS firewall, or edge gateway.

The model consumes pre-extracted heuristic indicators rather than raw lexical URL strings. A production implementation must include an upstream parser or feature extractor that obtains all 30 predictors in the same definitions, encoding, and order as the training data. External traffic reputation or domain-age signals may also require network lookups and can dominate the classifier's own inference time. Finally, the shared-covariance Gaussian assumption yields a linear boundary and may not represent nonlinear interactions or changes in phishing tactics; ongoing validation on newer, independently collected websites is needed before deployment.

\subsubsection{Figures}
The following figure environments use the requested placeholder paths. If an image is not yet present, the boxed fallback keeps this standalone document compilable.

\begin{figure}[!t]
\centering
\IfFileExists{figures/confusion_matrix.png}{\includegraphics[width=0.92\columnwidth]{figures/confusion_matrix.png}}{\fbox{\parbox[c][1.25in][c]{0.88\columnwidth}{\centering Placeholder: \texttt{confusion\_matrix.png}\\TN=1161, FP=71, FN=113, TP=866}}}
\caption{Held-out test confusion matrix for the selected LDA configuration.}
\label{fig:lda-confusion}
\end{figure}

\begin{figure}[!t]
\centering
\IfFileExists{figures/roc_curve.png}{\includegraphics[width=0.92\columnwidth]{figures/roc_curve.png}}{\fbox{\parbox[c][1.25in][c]{0.88\columnwidth}{\centering Placeholder: \texttt{roc\_curve.png}\\ROC-AUC = 0.9738}}}
\caption{Receiver Operating Characteristic curve (reported AUC 0.9738).}
\label{fig:lda-roc}
\end{figure}

\begin{figure}[!t]
\centering
\IfFileExists{figures/precision_recall.png}{\includegraphics[width=0.92\columnwidth]{figures/precision_recall.png}}{\fbox{\parbox[c][1.25in][c]{0.88\columnwidth}{\centering Placeholder: \texttt{precision\_recall.png}\\PR-AUC = 0.9694}}}
\caption{Precision--Recall curve (reported area 0.9694).}
\label{fig:lda-pr}
\end{figure}

\begin{figure}[!t]
\centering
\IfFileExists{figures/lda_comparison.png}{\includegraphics[width=0.92\columnwidth]{figures/lda_comparison.png}}{\fbox{\parbox[c][1.25in][c]{0.88\columnwidth}{\centering Placeholder: \texttt{lda\_comparison.png}\\Five-fold stratified CV F1 comparison: SVD and LSQR.}}}
\caption{Five-fold stratified training cross-validation comparison of SVD and LSQR with automatic shrinkage.}
\label{fig:lda-solvers}
\end{figure}

\end{document}
